\subsection{The full principal Cartan at every level}\label{sec:divisorial-full-cartan}
This subsection computes the fixed field of the full principal nonsplit Cartan at every level by constructing invariant coordinates with controlled errors.

\begin{proposition}\label{prop:divisorial-full-cartan}
Use the notation and valuation of Proposition~\ref{prop:divisorial-last-cartan}, with $p$ odd and $n\ge2$. Thus
\[
 A=p^2+p+1,\qquad B=p^2,\qquad
 w(x)=A,\quad w(y)=B,\quad
 w(p)=(p^4-1)p^{2n-2},
\]
where $x=x_n$, $y=y_n$, and $F$ is the completed valued field. Let $X,Y$ be the initial forms of $x,y$, put $t=\operatorname{res}(x^B/y^A)$, and write
\[
 \mathcal O=\Z_p[J],\qquad
 C=(1+p\mathcal O)/(1+p^n\mathcal O),\qquad r=p^{n-1}.
\]
Then
\begin{equation}\label{eq:full-cartan-field}
 \mathcal G(F^C)=k(X^r,Y^r),\qquad
 e(F/F^C)=r,\qquad \kappa(F^C)=k(t^r).
\end{equation}
In particular $f(F/F^C)=r$, and the residue extension is purely inseparable.
At $n=4$, let $S=(1+p\Z_p)/(1+p^4\Z_p)$ be the scalar subgroup. Then
\begin{equation}\label{eq:full-cartan-scalar-relative}
 e(F^S/F^C)=p^2,\qquad f(F^S/F^C)=p,
\end{equation}
and the latter residue extension is purely inseparable.
\end{proposition}

\begin{proof}
We use $\mathcal G(K)=\operatorname{Frac}(\operatorname{gr}_wK)$ with the restriction of the same valuation $w$ on every intermediate field.

\paragraph{Numerical and fixed-field preliminaries.}
We first record the bounds needed to lift initial forms inside actual fixed fields. Set
\[
 D=A-B=p+1,\qquad E=w(p)=(p^4-1)p^{2n-2},
 \qquad C_j=(1+p^j\mathcal O)/(1+p^n\mathcal O).
\]
Here $C_n=1$, $|C_j|=p^{2(n-j)}$, and $C_j/C_{j+1}$ is the additive group $\mathcal O/p\mathcal O$. Its independent generators are represented by
\[
 \sigma_j=1+p^j,\qquad \tau_j=1+p^jJ.
\]
All these elements commute. Their induced actions on $F^{C_{j+1}}$ have order $p$, and their joint fixed field is $F^{C_j}$. These degrees follow from the faithful finite group action on $F$. Faithfulness persists in the completion because it holds on the dense fraction field. The fixed subfields are complete.

The following inequalities hold for every $p\ge3$:
\begin{equation}\label{eq:full-cartan-elementary}
 p^2D>A,\qquad
 (p-1)B-pD=p(p^2-2p-1)>0,\qquad p^4-1>pA.
\end{equation}
If $0\le m\le n-2$, $r_m=p^m$, $j=n-m-1$, and $Q=p^{2j}$, they imply
\begin{equation}\label{eq:full-cartan-power-bound}
 E>p r_m Q A,
\end{equation}
since $p r_m Q=p^{2n-m-1}\le p^{2n-1}$.

For positive-valued $z,h$ and a power $s$ of $p$, all intermediate binomial coefficients in $(z+h)^s-z^s$ are divisible by $p$. If $E>s\,w(h)$, their terms have value greater than $s\,w(h)$, so
\begin{equation}\label{eq:full-cartan-ordinary-power}
 \operatorname{in}\bigl((z+h)^s-z^s\bigr)
       =\operatorname{in}(h)^s.
\end{equation}
For a lower bound, if $w(h)\ge c>0$ and $E>sc$, the same expansion gives $w((z+h)^s-z^s)\ge sc$. These statements concern ordinary powers in the characteristic-zero field, not formal multiplication by $p$.

We shall repeatedly use Lemma~\ref{lem:divisorial-controlled-lift}. If an order-$p$ automorphism $g$ has positive displacement shift $d$ and first derivation $D_g$ on $\mathcal G(K)$, then
\[
 \mathcal G(K^{\langle g\rangle})\subseteq\ker D_g.
\]
When $[\mathcal G(K):\ker D_g]=p$, equation~\eqref{eq:divisorial-graded-inequality} and $[K:K^{\langle g\rangle}]=p$ force equality. Only after establishing that equality do we choose an invariant lift $z_0$ with
\begin{equation}\label{eq:full-cartan-control}
 w(z_0-z)\ge w(gz-z)-d
\end{equation}
for an initial form in $\ker D_g$. A fixed $z$ needs no correction. Displacement bounds on homogeneous rational generators extend to the complete field by successive leading-term approximation, as proved above. The coefficient field $k$ has fixed Teichm\"uller representatives in $W(k)$.

\paragraph{The invariant-coordinate induction.}
We keep a quantitative comparison with the original coordinate powers. For $0\le m\le n-1$, put
\[
 K_m=F^{C_{n-m}},\qquad r_m=p^m.
\]
We prove
\begin{equation}\label{eq:full-cartan-induction-field}
 \mathcal G(K_m)=k(X^{r_m},Y^{r_m}),
\end{equation}
and choose $u_m,v_m\in K_m$ with those initial forms such that, for $m\ge1$,
\begin{equation}\label{eq:full-cartan-induction-error}
 \begin{split}
 w(u_m-x^{r_m})&\ge H_m,\\
 w(v_m-y^{r_m})&\ge H_m,
 \end{split}
 \qquad H_m=r_m p^{2(n-m)}D.
\end{equation}
At $m=0$, take $u_0=x$, $v_0=y$. The assertion about the graded field is already known. The errors are zero. For $m\ge1$, $H_m>r_mA>r_mB$, so the stated error bound preserves the initial forms.

Suppose the assertions hold at $m\le n-2$. Write
\[
 \rho=r_m,\qquad j=n-m-1,\qquad Q=p^{2j},
 \qquad u=u_m,\quad v=v_m,\quad U=X^{\rho},\quad V=Y^{\rho}.
\]
When $m\ge1$, the previous error bound satisfies
\begin{equation}\label{eq:full-cartan-old-error}
 H_m=\rho p^2QD>\rho QA.
\end{equation}
Thus valuation preservation alone makes the displacement of each old error too large to affect any of the coordinate initial forms below.

\paragraph{The induced coordinate displacements.}
The torsion initial forms from Proposition~\ref{prop:divisorial-last-cartan} and the identity
\[
 G(T,Z)-T=Z\bigl(1+\text{terms in }(T,Z)\bigr)
\]
give initial differences $X^Q,Y^Q$ for $\sigma_j$ on $x,y$, and $Y^Q,\bar\epsilon X^Q$ for $\tau_j$, with $\bar\epsilon\in\F_p^*$. Equations~\eqref{eq:full-cartan-power-bound} and \eqref{eq:full-cartan-ordinary-power} give their ordinary $\rho$-th powers. Equation~\eqref{eq:full-cartan-old-error} then permits replacement by the actual invariant coordinates $u,v$. The table of initial differences is
\begin{equation}\label{eq:full-cartan-displacements}
 \begin{array}{c|cc}
       &u&v\\ \hline
 \sigma_j-1&U^Q&V^Q\\
 \tau_j-1&V^Q&\bar\epsilon U^Q .
 \end{array}
\end{equation}
The weights of the four entries are $\rho QA,\rho QB,\rho QB,\rho QA$, respectively.

For $\tau=\tau_j$ on $K_m$, the smallest relative shift is
\[
 d_\tau=\rho(QB-A)>0.
\]
The shift on $v$ is larger by $\rho(Q+1)D$. The first derivation is therefore
\begin{equation}\label{eq:full-cartan-cross-derivation}
 D_\tau=V^Q\frac{\partial}{\partial U}
       \quad\text{on }k(U,V).
\end{equation}
For $\sigma=\sigma_j$ on $K_m$, the smallest relative shift and derivation are
\begin{equation}\label{eq:full-cartan-scalar-shift}
 d_\sigma=\rho(Q-1)B>0,\qquad
 D_\sigma=V^Q\frac{\partial}{\partial V}.
\end{equation}
The bounds hold on the entire complete field by the approximation argument stated above.

\paragraph{First fix the cross generator.}
We construct coordinates in $L=K_m^{\langle\tau\rangle}$ before treating the scalar action there. The kernel of \eqref{eq:full-cartan-cross-derivation} is $k(U^p,V)$, of index $p$ in $k(U,V)$. Since $[K_m:L]=p$, the graded-degree argument proves
\begin{equation}\label{eq:full-cartan-cross-field}
 \mathcal G(L)=k(U^p,V).
\end{equation}
The ordinary-power estimate gives $w(\tau(u^p)-u^p)=p\rho QB$, while the table gives $w(\tau v-v)=\rho QA$. Formula~\eqref{eq:full-cartan-control} produces $a,b\in L$ satisfying
\begin{equation}\label{eq:full-cartan-first-errors}
 \begin{split}
 w(a-u^p)&\ge E_a:=\rho((p-1)QB+A),\\
 w(b-v)&\ge E_b:=\rho(QD+A).
 \end{split}
\end{equation}
Here $E_a>p\rho A$ and $E_b>\rho B$. These strict inequalities verify the hypotheses of the controlled-lifting lemma and give initial forms $U^p,V$ and weights $p\rho A,\rho B$.

\paragraph{Then fix the scalar generator.}
We determine its action on the actual coordinates $a,b$. The Cartan is abelian, so $\sigma$ preserves $L$, acts there with order $p$, and has fixed field $K_{m+1}$. The bound on $K_m$ gives
\[
 w\bigl(\sigma(b-v)-(b-v)\bigr)\ge E_b+d_\sigma>\rho QB,
\]
where the strict margin is $E_b-\rho B=\rho(Q+1)D$. Hence
\begin{equation}\label{eq:full-cartan-scalar-b}
 \operatorname{in}(\sigma b-b)=V^Q.
\end{equation}
For the other coordinate, the ordinary-power estimate and the same displacement bound give
\begin{equation}\label{eq:full-cartan-scalar-a}
 w(\sigma a-a)\ge\min(p\rho QA,E_a+d_\sigma).
\end{equation}
Both terms on the right exceed $p\rho A+d_\sigma$, because
\[
 \begin{split}
 p\rho QA-p\rho A-d_\sigma&=\rho(Q-1)(pA-B)>0,\\
 E_a-p\rho A&=\rho(p-1)(QB-A)>0.
 \end{split}
\]
Thus on $\mathcal G(L)=k(U^p,V)$ the first scalar derivation still has shift $d_\sigma$, vanishes on $U^p$, and sends $V$ to $V^Q$. Its kernel is $k(U^p,V^p)$, of index $p$. The second actual degree-$p$ fixed-field extension therefore gives
\begin{equation}\label{eq:full-cartan-next-field}
 \mathcal G(K_{m+1})=k(U^p,V^p)=k(X^{p\rho},Y^{p\rho}).
\end{equation}

\paragraph{Errors for the next congruence layer.}
We complete the induction by constructing sufficiently accurate coordinates in $K_{m+1}$. Put $H=p\rho QD=H_{m+1}$. Applying the controlled-lifting lemma to $a$ inside $L$ gives an invariant $u'$ with
\begin{equation}\label{eq:full-cartan-second-u}
 w(u'-a)\ge\min(p\rho QA-d_\sigma,E_a).
\end{equation}
For $b^p$, equations~\eqref{eq:full-cartan-power-bound}, \eqref{eq:full-cartan-ordinary-power} and \eqref{eq:full-cartan-scalar-b} give $w(\sigma(b^p)-b^p)=p\rho QB$. An invariant $v'$ satisfies
\begin{equation}\label{eq:full-cartan-second-v}
 w(v'-b^p)\ge E_v:=p\rho QB-d_\sigma
                  =\rho((p-1)Q+1)B.
\end{equation}
The required error margins are
\begin{equation}\label{eq:full-cartan-error-margins}
 \begin{split}
 E_a-H&=\rho\{Q((p-1)B-pD)+A\}>0,\\
 p\rho QA-d_\sigma-H&=\rho\{(p-1)Q+1\}B>0,\\
 E_v-H&=\rho\{Q((p-1)B-pD)+B\}>0,\\
 pE_b-H&=p\rho A>0.
 \end{split}
\end{equation}
The first and third inequalities use \eqref{eq:full-cartan-elementary}. Since $E>p\rho QA>H$, the binomial expansion of $b^p-v^p$ gives
\[
 w(b^p-v^p)\ge\min(pE_b,E)>H.
\]
Equations~\eqref{eq:full-cartan-first-errors}, \eqref{eq:full-cartan-second-u} and \eqref{eq:full-cartan-error-margins} give $w(u'-u^p)>H$. Equations~\eqref{eq:full-cartan-second-v} and \eqref{eq:full-cartan-error-margins} similarly give $w(v'-v^p)>H$.

For $m\ge1$, taking ordinary $p$-th powers of the old approximation errors gives valuation at least $\min(pH_m,E)>H$, because $pH_m=p^2H$. At $m=0$ those errors vanish. Consequently
\[
 w(u'-x^{p\rho})\ge H,\qquad w(v'-y^{p\rho})\ge H.
\]
Taking $u_{m+1}=u'$ and $v_{m+1}=v'$ completes \eqref{eq:full-cartan-induction-field}--\eqref{eq:full-cartan-induction-error}.

\paragraph{Ramification and residue fields.}
The final induction step has $m=n-1$, so $K_m=F^C$ and $r_m=r$. Its graded field is $k(X^r,Y^r)$. The value group is generated by $rA,rB$ and is therefore $r\Z$, since $\gcd(A,B)=1$. Its homogeneous degree-zero field is
\[
 k\bigl((X^r)^B/(Y^r)^A\bigr)=k(t^r).
\]
This proves \eqref{eq:full-cartan-field}: the residue extension $k(t)/k(t^r)$ is purely inseparable of degree $r$. The product $ef=r^2$ agrees with the group order after the fixed field has been computed.

At $n=4$, Proposition~\ref{prop:divisorial-scalar-level-four} gives $e(F/F^S)=p$ and $f(F/F^S)=p^2$. Multiplicativity in the actual tower $F/F^S/F^C$ now yields \eqref{eq:full-cartan-scalar-relative}. The intermediate residue extension is $k(t^{p^2})/k(t^{p^3})$, hence purely inseparable.
\end{proof}

The exact Cartan ramification degree gives a sharper conditional index for the further tame quotient.
\begin{corollary}\label{cor:divisorial-full-cartan-6500}
Take $p=5$ and $n=4$ in Proposition~\ref{prop:divisorial-full-cartan}. If $F^C/K$ is a further tame extension of ramification index $12$ with unchanged residue field, then the absolute ramification index of $K$ is exactly $6500$.
\end{corollary}
\begin{proof}
The valuation on $F$ gives $w(5)=624\cdot5^6$, and the proposition gives $e(F/F^C)=5^3$. Thus
\[
 e_K=\frac{624\cdot5^6}{12\cdot5^3}=6500.
\]
The existence of the further tame extension is a hypothesis. This calculation concerns the specified divisorial valuation and supplies no perfect-residue specialization or actual unequal-character boundary pair.
\end{proof}
